\section{Arquitetura do Sistema e Integração entre Serviços}
\label{sec:2.4}

O desenvolvimento de sistemas de software complexos demanda escolhas arquiteturais capazes de equilibrar organização interna do código, desempenho operacional, manutenibilidade e capacidade de evolução. No caso do IGRIS, essa estrutura é formada pela combinação entre um monolito modular, responsável por concentrar os principais módulos da aplicação, e adaptadores de microsserviços, utilizados para integrar as diferentes implementações dos algoritmos genéticos.

O monolito modular é uma variação do monolito tradicional na qual a aplicação é executada como um único processo, mas organizada internamente em módulos separados e coesos. Cada módulo possui responsabilidades próprias e pode ser desenvolvido de forma relativamente independente, embora todos sejam combinados em um único artefato para implantação \cite[p.~15--16]{newman2021}. Diferentemente do monolito clássico, em que o crescimento do código tende a formar uma estrutura homogênea e fortemente acoplada, o monolito modular estabelece fronteiras explícitas entre responsabilidades de negócio dentro do mesmo processo. Essa organização preserva a clareza estrutural e reduz a complexidade operacional associada a arquiteturas distribuídas completas.

A arquitetura em camadas, conhecida como \textit{Layered Architecture}, organiza o fluxo de dados dentro dos módulos do sistema. Nessa abordagem, as responsabilidades são distribuídas entre camadas com papéis definidos: a camada de \textit{Controller} recebe e direciona as requisições HTTP; a camada de \textit{Service} concentra a lógica de negócio; a camada de \textit{Repository} abstrai o acesso ao banco de dados; e a camada de \textit{Model} representa as entidades do domínio \cite[Seção~6.3.1]{sommerville2016}. Essa separação reduz o acoplamento entre responsabilidades distintas, facilita a manutenção do código e favorece a realização de testes unitários.

\subsection{Adaptadores de Microsserviços e Comunicação REST}
\label{subsec:2.4.1}

Quando parte do processamento apresenta custo computacional elevado, como ocorre na geração de grades por Algoritmos Genéticos, sua execução pode comprometer a responsividade do sistema principal. Nesse cenário, o uso de serviços externos permite isolar esse processamento em componentes independentes, com ciclo de vida, linguagem e ambiente de execução próprios \cite[p.~3]{newman2021}. No IGRIS, essa integração ocorre por meio de adaptadores, componentes responsáveis por encapsular os detalhes de comunicação HTTP e traduzir as interfaces dos serviços externos para o vocabulário interno da aplicação.

A comunicação entre o sistema principal e os serviços externos ocorre via HTTP, seguindo princípios associados ao estilo REST, ou \textit{Representational State Transfer}. Nesse modelo, as requisições devem conter as informações necessárias para seu processamento, reduzindo a dependência de estado armazenado no servidor. Além disso, os recursos e operações são acessados por meio dos métodos semânticos do protocolo HTTP, como GET, POST, PUT e DELETE \cite[p.~3]{newman2021}. No contexto do IGRIS, o sistema principal envia aos serviços algorítmicos um \textit{payload} estruturado em JSON contendo os dados necessários para a geração da grade, e recebe como resposta os resultados processados em formato igualmente estruturado.

A containerização por meio do Docker permite que serviços implementados em linguagens distintas operem em ambientes isolados e reprodutíveis, comunicando-se pela rede interna da aplicação. Essa estratégia é especialmente relevante para o IGRIS, pois possibilita executar implementações desenvolvidas em Java, Python e Go sem exigir que todas compartilhem o mesmo ambiente de dependências. Além disso, a utilização de containers contribui para padronizar as condições de execução nos experimentos de benchmark, reduzindo interferências externas e tornando mais consistente a comparação entre as diferentes implementações do algoritmo genético \cite[p.~20]{newman2021}.
